% Figure: oscillatory pipe velocity profiles at low and high Womersley number,
% sketched at one instant of the cycle. At low Wo the profile tracks the
% quasi-steady Poiseuille parabola (viscosity fills the bore); at high Wo the
% core moves as a flat-topped slug driven by inertia with a thin oscillatory
% Stokes layer at the wall. Plotted schematically as u (arbitrary, normalised)
% against radius r/R with the axis at 0 and walls at plus/minus 1.
\begin{figure}[htbp]
\centering
\begin{tikzpicture}
\begin{axis}[
  width=11cm, height=7cm,
  xlabel={radius $r/R$ (walls at $\pm1$, axis at $0$)},
  ylabel={velocity (instantaneous, normalised)},
  xmin=-1, xmax=1, ymin=0, ymax=1.15,
  legend pos=south east,
  legend cell align=left,
  grid=both, grid style={enggray!18},
  minor grid style={enggray!8},
  tick label style={font=\scriptsize},
  label style={font=\small},
  legend style={font=\scriptsize},
]
% Low Womersley: quasi-steady parabola u = 1 - (r/R)^2.
\addplot[engred, very thick, domain=-1:1, samples=80] {1 - x^2};
\addlegendentry{low $\mathrm{Wo}\ll1$ (parabolic)}
% High Womersley: flat slug core with thin wall layer. Approximate with a
% high-power profile 1 - (r/R)^{16}, flat over the core and steep at the wall.
\addplot[engblue, very thick, domain=-1:1, samples=160] {1 - x^16};
\addlegendentry{high $\mathrm{Wo}\gg1$ (slug + wall layer)}
% Mark the thin Stokes wall layer on the right.
\draw[enggray, thick, {Latex}-{Latex}] (axis cs:0.86,0.5) -- (axis cs:1.0,0.5);
\node[font=\scriptsize, enggray, anchor=south] at (axis cs:0.9,0.52) {$\delta_s$};
\end{axis}
\end{tikzpicture}
\caption[Pulsatile pipe profiles versus Womersley number]{Instantaneous
oscillatory pipe profiles at low and high Womersley number $\mathrm{Wo} =
R\sqrt{\omega/\nu}$. At low $\mathrm{Wo}$ viscosity diffuses momentum across
the bore within a fraction of a cycle, so the profile tracks the quasi-steady
Poiseuille parabola at every phase. At high $\mathrm{Wo}$ the oscillation
outruns viscous diffusion, and the core moves as a flat-topped inertial slug
that leads the pressure gradient in phase, with the oscillatory shear confined
to a thin Stokes layer $\delta_s = \sqrt{2\nu/\omega}$ at the wall. Worked
example~23 places the aorta at $\mathrm{Wo}\approx19$ and a microvessel at
$\mathrm{Wo}\approx0.08$.}
\label{fig:vol03ch12:womersley-profiles}
\end{figure}
